\section{Типы данных и их размеры}

Базовые типы данных:
\begin{description}
\item[\texttt{char}] --- один байт, содержащий один символ из локального символьного набора.
\item[\texttt{int}] --- целое число, обычно имеющее типовой размер для целых чисел в данной системе.
\item[\texttt{float}] --- вещественное число одинарной точности с плавающей точкой.
\item[\texttt{double}] --- вещественное число двойной точности с плавающей точкой.
\end{description}

Модификаторы \texttt{short} и \texttt{long} могут применяться к типу \texttt{int}
для разграничения длины целых чисел (слово \texttt{int} в таких объявлениях обычно опускается).
Компилятор самостоятельно выбирает размер типа в соответствии с характеристиками аппаратуры
и следующими ограничениями: \texttt{short} и \texttt{int} должны иметь длину не менее 16 бит,
\texttt{long} --- не менее 32 бит; \texttt{short} должен быть не больше \texttt{int},
a \texttt{int} --- не больше \texttt{long}.
Обычно \texttt{short} занимает 16 бит, \texttt{int} --- 16 или 32 бита, \texttt{long} --- 32 бита. 

Модификатор \texttt{signed} (<<со знаком>>) или \texttt{unsigned} (<<без знака>>) может 
применяться к типу \texttt{char} или любому целочисленному.
Числа типа \texttt{unsigned} всегда неотрицательны, а длина диапазона их значений
равна степени двойки $2^n$, где $n$ --- количество битов в машинном представлении типа.

Тип \texttt{long double} обозначает число с плавающей точкой повышенной точности.
Как и в случае целочисленных переменных, длины вещественных объектов
зависят от реализации языка, так что из длин типов \texttt{float}, \texttt{double}
и \texttt{long double} различными могут быть одна, две или три.

Стандартные заголовочные файлы \texttt{<limits.h>} и \texttt{<float.h>}
содержат символические константы для всех этих размеров, а также других свойств
системы и компилятора.

